\documentclass{article}
\usepackage[utf8]{inputenc}
\usepackage{geometry}
\geometry{a4paper, margin=1in}
\usepackage{minted}

\title{Chapter 1 Summary}
\author{Davide Colabove}
\date{}

\begin{document}

\maketitle

\section*{UAV Software Engineering: Chapter 1 Overview}
This document serves as a brief summary of the key concepts covered in Chapter 1, designed for quick reference on the portfolio website.

\section{Memory Management and RAII}
In critical UAV systems, memory leaks are unacceptable. Modern C++ employs the \textbf{RAII (Resource Acquisition Is Initialization)} paradigm. Memory allocation occurs in the constructor, and deallocation happens in the destructor. When an object goes out of scope, memory is automatically freed, ensuring consistent execution times for high-frequency control loops.[cite: 3]

\subsection*{Smart Pointers}
We explored \texttt{std::unique\_ptr}, which grants exclusive ownership of a heap-allocated memory block without adding overhead compared to raw pointers. Because ownership is exclusive, it cannot be copied, only moved via \texttt{std::move()}, preventing dangling references and double-free errors.[cite: 3] 
Modern C++ best practices dictate using \texttt{std::make\_unique} for exception safety. Other system components can read data safely without taking ownership using the \texttt{get()} method.[cite: 3]

\section{Concurrency and Thread Safety}
Multicore environments introduce the risk of \textbf{data races}, where multiple threads access the same memory location simultaneously (with at least one write operation). In a flight controller, this can lead to corrupted data (like altitude readings) and erratic motor commands.[cite: 1]

\subsection*{Synchronization Primitives}
To ensure thread safety, C++ provides several tools:
\begin{itemize}
    \item \textbf{\texttt{std::mutex}:} Ensures mutual exclusion for critical sections. To avoid deadlocks (where a lock is acquired but never released), we use \texttt{std::lock\_guard<std::mutex>} to apply RAII to mutex locking, guaranteeing exception-safe unlocking.[cite: 1]
    \item \textbf{\texttt{std::atomic}:} For lock-free, hardware-level atomic operations on simple variables.[cite: 1]
    \item \textbf{\texttt{std::condition\_variable}:} Used to block a thread until a specific event occurs, often paired with a mutex.[cite: 1]
\end{itemize}

\subsection*{Producer-Consumer Architecture}
To handle asynchronous data generation (e.g., from sensors) and processing, we implemented the \textbf{Producer-Consumer pattern} using a thread-safe \texttt{std::queue} (FIFO container) and a \texttt{std::condition\_variable}. This ensures telemetry packets are buffered and processed sequentially without data loss or wasted CPU cycles.[cite: 1]

\section{Build Systems: CMake}
Managing large C++ projects (like ROS 2 workspaces) requires robust build systems. \textbf{CMake} is the industry standard—a cross-platform generator that reads high-level rules (\texttt{CMakeLists.txt}) and produces low-level build files (like Makefiles).[cite: 2]

\subsection*{Out-of-Source Builds}
To keep the source code directory clean, we enforce \textbf{out-of-source builds}:
\begin{enumerate}
    \item Create an isolated directory: \texttt{mkdir build \&\& cd build}.[cite: 2]
    \item Generate build files targeting the parent directory: \texttt{cmake ..}.[cite: 2]
    \item Compile the binary: \texttt{make}.[cite: 2]
\end{enumerate}
The \texttt{build/} directory should always be added to \texttt{.gitignore} to prevent compiled binaries from polluting the version control history.[cite: 2]

\end{document}